\chapter{Genetic Scheduling Algorithm}
\label{ch:method}

\section{Overview}
The algorithm evolves a population of expression trees. Each generation:
(1) take a snapshot of the world (satellites and tasks); (2) decode a schedule
with every tree and compute its fitness; (3) breed the next population.

\section{Individual Representation}
An individual is one tree $T$. Internal nodes are operators
$\{+,-,\times,\div,\max,\min\}$ (binary) and $\{\text{neg},|\cdot|,\text{inv}\}$
(unary); leaves are terminals read from a \emph{registry}.
Ephemeral constants are disabled. The terminals are:
\begin{itemize}
  \item task: arrival time, RAM, disk, image size, execution time, transfer time, number of hops;
  \item satellite: elevation angle, latency, bandwidth, completed tasks, CPU/network busy-until,
        CPU/network capacity, reserved and remaining energy, load ratio.
\end{itemize}
The tree evaluated on $(j,s)$ returns a score: a higher score means the pair is preferred.

\section{Greedy Decoder}
\begin{algorithm}[h]
\caption{Greedy decoding of a tree $T$}
\begin{algorithmic}[1]
\State mark tasks arrived after the snapshot time as unavailable
\State compute $score[j][s] \gets T(j,s)$ for all available $j$, all $s$
\While{unscheduled available tasks remain}
  \State among feasible pairs $(j,s)$ pick the one with highest $score[j][s]$
  \If{none} \textbf{break} \EndIf
  \State assign $j$ to $s$; update $f_s$, $E_s$, completed tasks of $s$
  \State recompute column $score[\cdot][s]$ for unscheduled tasks
\EndWhile
\end{algorithmic}
\end{algorithm}
Only the committed satellite changes, so only its column of the score matrix is
refreshed. Tree evaluations drop from $O(J^2 S)$ to $O(JS + J^2)$ for $J$ tasks
and $S$ satellites, with identical results.

\subsection{Selection rule and ties}
\label{sec:ties}
Let $\mathcal{F}$ be the set of feasible pairs of a round, tasks and satellites
being numbered in the order of the input. ``Highest score'' in line~4 is made
precise by a lexicographic maximum. The \emph{default} rule is
\begin{equation}
  (j^\star,s^\star)=\operatorname*{lex\,max}_{(j,s)\in\mathcal{F}}
  \bigl(T(j,s),\,-j,\,-s\bigr),
  \label{eq:sel-index}
\end{equation}
so a tie goes to the lowest task index and then to the lowest satellite index.
(In the implementation a pair must beat $-\infty$ strictly, so a pair whose
score is $-\infty$ or not a number is never selected.) All experiments of
Chapter~\ref{ch:exp} use this rule.

A tree that does not read any satellite terminal gives every satellite the same
score. Under \eqref{eq:sel-index} it sends each task to the first satellite that
is still feasible, while the others stay idle: the satellite is chosen by the
input order, not by the tree. The optional \emph{earliest-free} rule removes this
artefact:
\begin{equation}
  (j^\star,s^\star)=\operatorname*{lex\,max}_{(j,s)\in\mathcal{F}}
  \bigl(\bar T(j,s),\,-f_s,\,-j,\,-s\bigr),
  \label{eq:sel-free}
\end{equation}
where $f_s$ is the time at which $s$ becomes free and $\bar T$ equals $T$ with
``not a number'' replaced by $-\infty$, so that the order is total and a pair is
selected whenever $\mathcal{F}\neq\emptyset$. Among equally scored pairs the one
that can start first wins. A strictly higher score always wins, so the tree
keeps full control over every choice it actually expresses.

Sending a task to the machine that becomes available first is the machine
choice of list scheduling \citep{graham1966bounds}. Since $d_j$ does not depend
on the satellite, earliest start and earliest finish coincide, and the rule is
also the processor selection of the earliest-finish-time and
minimum-completion-time heuristics
\citep{topcuoglu2002heft,braun2001comparison}. Using it only to break ties, and
the particular key \eqref{eq:sel-free}, are choices of this work.

\paragraph{Proposition 2 (task-only rules without constraints).} Assume that no
pair is ever infeasible ($\omega_s=E_s=\infty$) and that the score is a finite
function of the task only, $T(j,s)=g(j)$. Then both rules schedule the tasks in
the same order $\pi$ (by decreasing $g$, ties by index), and for every $k$
\begin{equation}
  c^{\mathrm{free}}_{\pi(k)}\;\le\;t_0+\frac1S\sum_{i<k}d_{\pi(i)}+d_{\pi(k)}
  \;\le\;t_0+\sum_{i\le k}d_{\pi(i)}\;=\;c^{\mathrm{index}}_{\pi(k)} .
\end{equation}
Every task completes no later under \eqref{eq:sel-free}; hence $R$ is not larger,
and $C=1$ under both rules.

\emph{Proof.} With all pairs feasible, the pairs of maximal score are those of the
tasks with the largest $g$, each combined with every satellite. Under both rules
the task selected is the one of lowest index among them, which gives the same
order $\pi$. Under \eqref{eq:sel-index} the satellite is always the first one, so
the tasks run in sequence on it and $c^{\mathrm{index}}_{\pi(k)}=t_0+\sum_{i\le k}d_{\pi(i)}$.
Under \eqref{eq:sel-free} the $k$-th task starts at $\min_s f_s$. The satellites
are never idle before their last task, so $\sum_s(f_s-t_0)=\sum_{i<k}d_{\pi(i)}$,
and the minimum is at most the mean: $\min_sf_s\le t_0+\frac1S\sum_{i<k}d_{\pi(i)}$.
\hfill$\square$

The schedule of \eqref{eq:sel-free} is then exactly the list schedule of the list
$\pi$, whose makespan is at most $2-1/S$ times the optimum
\citep{graham1966bounds}.

\paragraph{No guarantee with constraints.} With the energy filter, spreading the
tasks can schedule \emph{fewer} of them. Take two satellites with $E_s=30$,
$\rho_s=1$, and three tasks of duration $10$, $10$, $30$ under shortest-task-first.
Rule \eqref{eq:sel-index} puts both short tasks on the first satellite and the
long one fits on the second ($C=1$). Rule \eqref{eq:sel-free} puts one short task
on each satellite, and no satellite has $30$ units left ($C=2/3$). Filling one
satellite first keeps the budget of the other whole, as first-fit does in bin
packing. In the data of Chapter~\ref{ch:exp} the energy budget is the binding
constraint, so the earliest-free rule is expected to lower response times and
may lower the completion ratio. Its effect has not been measured; it is off by
default. The bounds of Section~\ref{sec:fixed} hold under both rules, since
Proposition~1 does not depend on which feasible pair is selected.

\section{Fitness}
For individual $i$ with raw objectives $(C_i, R_i, \mathcal{E}_i)$, the response
time and energy are normalised across the population and
\begin{equation}
  f_i = w_C\, C_i - w_R\, \hat R_i - w_E\, \hat{\mathcal{E}}_i ,
\end{equation}
with default weights $w_C=0.60$, $w_R=0.25$, $w_E=0.15$.
Here $\hat R_i$ and $\hat{\mathcal{E}}_i$ are min--max normalised over the current
population $P$, e.g.\ $\hat R_i = (R_i-\min_{k\in P}R_k)/(\max_{k\in P}R_k-\min_{k\in P}R_k)$;
an individual that schedules nothing gets $\hat R_i=1$. We call this the
\emph{population-relative} fitness. It is the one used in all experiments of
Chapter~\ref{ch:exp}.

\subsection{Fixed-bound fitness}
\label{sec:fixed}
The population-relative fitness has two weaknesses, both visible in
Chapter~\ref{ch:exp}. (i) $R$ is a mean over the scheduled tasks only, so a rule
that schedules few, short tasks is rewarded. (ii) Min--max normalisation maps
any spread of a cost to $[0,1]$, however small: $\mathcal{E}$ differs by about
$1\%$ between rules, yet it moves the fitness by up to $w_E$, more than any
realistic difference in $C$. Moreover $f_i$ depends on the other individuals, so
its values cannot be compared between generations.

The optional \emph{fixed-bound} fitness replaces the population statistics by
bounds that depend on the world only. This is the standard normalisation for
weighted sums, where each objective is divided by the distance between a lower
and an upper bound of its range \citep{marler2004survey}, and it is how
\citet{icdcs_sc} combine time and energy in the same setting (time over the
largest response, energy over the battery). Tasks that are not scheduled are
charged a fixed penalty that is added to the time criterion, as in scheduling
with rejection \citep{shabtay2013rejection}. The two bounds below are derived
here from the constraints of the decoder; they are not taken from those papers.

\paragraph{Bounds.} Let $t_0$ be the snapshot time, $\mathcal{J}=\{j: a_j\le t_0\}$
the available tasks, $D=\sum_{j\in\mathcal{J}}d_j$ and $a_{\min}=\min_{j\in\mathcal{J}}a_j$.
For each satellite let
\begin{equation}
  T_s=\max\!\Bigl(0,\;\min\bigl(D,\;\omega_s-t_0,\;E_s/\rho_s\bigr)\Bigr),
  \label{eq:Ts}
\end{equation}
where the last term is dropped if $\rho_s\le0$, and define
\begin{equation}
  R^{\max}=(t_0-a_{\min})+\max_s T_s,\qquad
  \mathcal{E}^{\max}=\min\Bigl(\sum_s\rho_sT_s,\;\rho_{\max}D\Bigr).
  \label{eq:bounds}
\end{equation}

\paragraph{Proposition 1 (validity).} For every tree, the schedule produced by
the greedy decoder satisfies $c_j-a_j\le R^{\max}$ for each scheduled task and
$\mathcal{E}\le\mathcal{E}^{\max}$.

\emph{Proof.} The decoder sets $f_s=t_0$ for every satellite and starts a task
at $f_s$, so the tasks of $s$ run back to back and a task on $s$ completes at
$t_0+L_s$, where $L_s$ is the total duration committed to $s$ up to and
including that task. Three facts bound $L_s$: only available tasks are
scheduled, so $L_s\le D$; the window constraint gives $t_0+L_s\le\omega_s$; and
the energy constraint is checked against the \emph{remaining} energy, so the
energy committed to $s$ never exceeds its initial budget, $\rho_sL_s\le E_s$.
Hence $L_s\le T_s$ and $c_j\le t_0+\max_sT_s$; with $a_j\ge a_{\min}$ this gives
the first claim. For the second, $\mathcal{E}=\sum_s\rho_sL_s\le\sum_s\rho_sT_s$,
and also $\mathcal{E}\le\rho_{\max}\sum_sL_s\le\rho_{\max}D$ because each task
runs on one satellite. \hfill$\square$

\paragraph{Fitness.} With $N$ tasks in the world and $\mathcal{A}_i$ the schedule
of individual $i$,
\begin{align}
  \tilde R_i &= \frac1N\Bigl[\sum_{(j,s)\in\mathcal{A}_i}\frac{c_j-a_j}{R^{\max}}
                 +\bigl(N-|\mathcal{A}_i|\bigr)\Bigr]
              = 1-C_i\Bigl(1-\frac{R_i}{R^{\max}}\Bigr), \label{eq:Rtilde}\\
  \tilde{\mathcal{E}}_i &= \mathcal{E}_i/\mathcal{E}^{\max},\qquad
  f^{\mathrm{fix}}_i = w_C\,C_i-w_R\,\tilde R_i-w_E\,\tilde{\mathcal{E}}_i . \label{eq:ffix}
\end{align}
An unscheduled task contributes $1$ to the bracket, i.e.\ it is charged the
response $R^{\max}$. With several worlds $w$ the bracket and $N$ are summed over
the worlds, each response being divided by the bound $R^{\max}_w$ of its own
world, and $\tilde{\mathcal{E}}_i=\sum_w\mathcal{E}_{i,w}/\sum_w\mathcal{E}^{\max}_w$.

\paragraph{Properties.}
\begin{enumerate}
  \item \emph{Range.} By Proposition~1, $\tilde R_i,\tilde{\mathcal{E}}_i\in[0,1]$,
        so $f^{\mathrm{fix}}_i\in[-(w_R+w_E),\,w_C]$. The empty schedule has
        $\tilde R=1$, the worst value, without a special case.
  \item \emph{No reward for dropping tasks.} A scheduled task costs
        $(c_j-a_j)/R^{\max}\le1$ in \eqref{eq:Rtilde}, an unscheduled one costs
        $1$. Equivalently, for a fixed mean response $R_i$, $\tilde R_i$
        decreases when $C_i$ grows. This removes weakness~(i).
  \item \emph{Independence of the population.} $f^{\mathrm{fix}}_i$ is a function
        of tree $i$ and of the worlds only. Hence the order of two individuals
        cannot be changed by a third one. Under the population-relative fitness
        it can: $f_a-f_b=w_C\Delta C-w_R\Delta R/\mathrm{range}(R)-w_E\Delta\mathcal{E}/\mathrm{range}(\mathcal{E})$,
        and an outlier that widens one range but not the other can change the
        sign.
  \item \emph{Monotone best fitness.} With at least one elite and an unchanged
        set of worlds, $\max_i f^{\mathrm{fix}}_i$ is non-decreasing over the
        generations, because the elite is copied unchanged and re-evaluated to
        the same value. The stagnation test of Section~\ref{sec:worlds} is then
        exact.
  \item \emph{Scale.} A relative change of $\mathcal{E}$ by $\varepsilon$ changes
        the fitness by at most $w_E\varepsilon$, instead of up to $w_E$. This
        removes weakness~(ii).
\end{enumerate}
The fitness is still a weighted sum: it can only reach the convex part of the
Pareto front \citep{das1997weighted}, and the weights remain a modelling choice.
Tournament selection only uses the order of the fitness values, so the change
affects selection through the ranking alone. The effect on the test metrics
has not been measured yet; the option is off by default.

\section{Evolution}
\begin{description}
  \item[Initialisation] ramped half-and-half, depths from 2 to the maximum depth.
  \item[Elitism] the top 2 individuals are copied unchanged.
  \item[Selection] tournament of size 5.
  \item[Crossover] subtree swap, applied with probability 0.80.
  \item[Mutation] per-node probability 0.08: terminals are redrawn; operators are
        swapped for another of the same arity, or (10\% of cases) the subtree is
        replaced by a fresh \emph{grow} tree of depth 3.
  \item[Hoist] with probability 0.05 per offspring, to fight bloat.
\end{description}
Every offspring satisfies $\text{depth}\le D_{\max}$ and
$\text{size}\le 2^{D_{\max}+1}$; a child that violates a cap is replaced by a copy
of its parent. Trees are algebraically simplified after mutation to remove redundancy.

\section{Changing Worlds and Convergence}
\label{sec:worlds}
The population is evaluated on a snapshot of one world. In \texttt{main.py}, if the best fitness
in the current world does not improve for a number of generations (the \emph{convergence
threshold}, default 10), the next pre-processed world is loaded. The population therefore faces
different scenarios over time, which should limit overfitting to one instance. Because fitness is
normalised within the population, it is comparable only inside one world; for this reason the
experiments of Chapter~\ref{ch:exp} use a fixed switching schedule and held-out test worlds.
With the fixed-bound fitness of Section~\ref{sec:fixed} the values are comparable between
generations of the same world (but still not between worlds, whose bounds differ).
